\section{Experiments}
\label{sec:experiments}

We compare G34 with direct recovery from within-neuron weight ratios
(Pair recovery). Both receive the same fitted Dense models; neither
observes a true indexator. A target task receives only the selected
indexator. Its codebook, hidden biases, output matrix and offset are
initialized and fitted afresh. Dense is the unstructured reference.

On images, both extractors consider 8, 16, 32 and 64 independent codes
and retain at most eight candidates per budget. Each candidate receives
18 numerical starts and 2000 Adam updates. Separate source samples
fit numerical parameters, select their checkpoints and select the
indexator. Development episodes choose the code budget and one global
learning rate, activation quantile and update count, averaging two
fresh starts. Test episodes use this fixed rule and one fresh start;
their query responses are opened after all predictions are frozen.
G34 calibrates the aligned source models jointly; Pair recovery builds
proposals separately from two source models ranked by validation loss.
Tables count all fitted numerical parameters, including biases and
the output layer. Image accuracy is reported with 95\% episode intervals,
conditional on the fixed source cohort and class split; images and
class pairs may repeat. Source validation assesses predictive fitting,
while exact weight identification and global convergence are not established.

\subsection{Synthetic recovery}

We use two fixed random sharing geometries satisfying the pair-identification
conditions, with 16 inputs, 16 hidden
ReLU units and four independent codes, with one and two outputs,
respectively. Each geometry supplies one new numerical source function
and 4096 fitting, 512 numerical-selection and 1024 structure-selection
examples. Dense preparation uses 12 Adam starts, 2500 updates,
Gauss--Newton refinement and two repair rounds. G34 and Pair recovery
use the same descriptors and at most 32 candidate indexators; each
candidate receives four numerical starts and 48 Gauss--Newton updates.
Only dimensions and the code budget are given to the learners.
Table~\ref{tab:recovery-synthetic} reports the subsequent structural audit,
allowing hidden-unit permutations and code-pair renaming and sign changes.
Both methods recover the full signed indexator, including its zero entries.
These are two geometric cases, not a statistical sample of all structures.

\begin{table}[htbp]
    \centering\small
    \begin{tabular}{lrrcc}
        \toprule
        Geometry & Active weights & Dense error & G34 & Pair recovery\\
        \midrule
        A & 160 & $1.65\times10^{-7}$ & Exact & Exact\\
        B & 135 & $9.99\times10^{-12}$ & Exact & Exact\\
        \bottomrule
    \end{tabular}
    \caption{Recovery from one fitted source function per geometry.
        Dense error is the maximum log-magnitude discrepancy after
        alignment, measured against generating weights only in this
        late audit. Pair recovery produces three and one complete
        candidates; G34 numerically compares 32 in each case.}
    \label{tab:recovery-synthetic}
\end{table}

\FloatBarrier
\subsection{MNIST class-pair transfer}

We use original $28\times28$ MNIST images~\cite{lecun1998gradient}
and a 784--16--2 ReLU network. Classes 0--4 define ten source pairs,
classes 5--6 define development episodes, and classes 7--9 define test
pairs. Each source pair has 8000 fitting examples and two disjoint
1000-example samples for numerical and structure selection. Dense
preparation compares 12 Adam configurations over 2500 updates and
continues four retained fits with up to 200 full-gradient L-BFGS iterations,
selecting numerical states by the source validation loss.
Sixteen development episodes and 96 test episodes use 16 training
images and 100 query images per class. Test training images come
from the official training split and query images from the official
test split. This is a fully connected, class-disjoint 2-way 16-shot
comparison; published standard MNIST scores are not matched baselines.
The shared source Dense preparation reaches 99.67\% mean numerical-validation accuracy across the ten class pairs (minimum 99.30\%), with mean cross-entropy 0.0109.

\begin{table}[htbp]
\centering\small
\begin{tabular}{lrrrrr}
\toprule
Method & Codes & Parameters & Updates & Accuracy (\%) & CE\\
\midrule
Dense & --- & 12594 & 50 & $92.60\pm0.74$ & 0.208\\
G34 & 64 & 114 & 175 & $89.49\pm0.78$ & 0.284\\
Pair recovery & 64 & 114 & 100 & $86.51\pm1.24$ & 0.341\\
\bottomrule\end{tabular}
\caption{MNIST: query prediction after fresh numerical fitting on 96 held-out class-pair episodes. Pair recovery minus G34: -2.98 percentage points, paired 95\% interval [-4.21, -1.76].}
\label{tab:recovery-mnist}
\end{table}


\FloatBarrier
\subsection{Fashion-MNIST class-pair transfer}

Fashion-MNIST~\cite{xiao2017fashion} supplies a second recognition task
with the same image size, architecture, class split, sample counts and
selection budget. Source classes are T-shirts, trousers, pullovers,
dresses and coats; development uses sandals and shirts; test pairs
use sneakers, bags and ankle boots. Thus this comparison transfers
constraints across held-out clothing classes, with all numerical
parameters learned anew. The two datasets use separately fitted
Dense models and separately selected structures and learning rules.
The shared source Dense preparation reaches 97.21\% mean numerical-validation accuracy across the ten class pairs (minimum 89.00\%), with mean cross-entropy 0.0744.

\begin{table}[htbp]
\centering\small
\begin{tabular}{lrrrrr}
\toprule
Method & Codes & Parameters & Updates & Accuracy (\%) & CE\\
\midrule
Dense & --- & 12594 & 100 & $96.40\pm0.73$ & 0.135\\
G34 & 64 & 114 & 50 & $94.45\pm0.73$ & 0.200\\
Pair recovery & 64 & 114 & 225 & $92.74\pm1.00$ & 0.211\\
\bottomrule\end{tabular}
\caption{Fashion-MNIST: query prediction after fresh numerical fitting on 96 held-out class-pair episodes. Pair recovery minus G34: -1.70 percentage points, paired 95\% interval [-2.53, -0.88].}
\label{tab:recovery-fashion}
\end{table}

\begin{figure}[htbp]
    \centering
    \includegraphics[width=\linewidth]{recovery_dense_quality.pdf}
    \caption{Source Dense fitting on the two image tasks. Each curve
        averages the ten class pairs using each pair's best Adam
        configuration by source validation. Both extractors subsequently
        receive the same validation-retained full-gradient continuations.
        These are source losses, not test results.}
    \label{fig:recovery-dense-quality}
\end{figure}

